Die Zusammenarbeit im Projektteam war stets konstruktiv und zielführend. Unterschiedliche Meinungen konnten sachlich diskutiert und gemeinsame Lösungen gefunden werden. Der regelmäßige Austausch in Online-Konferenzen war stets unkompliziert und zuverlässig. Das zeigt sich unter anderem daran, dass geplante Treffen nie verschoben werden mussten und am Ende jeder Konferenz der nächste Termin vereinbart wurde. Auch auf zwischenmenschlicher Ebene ließ sich gut zusammenarbeiten, da die Konferenzen auf mich stets locker und natürlich wirkten.

Gereon Such brachte seine Erfahrung in der Signalverarbeitung und im Entwurf von Leiterplattenlayouts gewinnbringend ein, da ohne sie einige Arbeiten deutlich länger gedauert hätten. Er übernahm die Auswahl und Integration der Sensoren und Treibermodule sowie die Erstellung des Schaltplans. Darüber hinaus setzte er sich früh mit dem Einsatz von KI zur Aufbereitung von Arbeitsergebnissen in unterschiedlichen Darstellungsformen auseinander. Damit half er dem Projektteam, zeitintensive Routinearbeiten schneller zu erledigen.

Stefan Etringer leistete gründliche Recherchearbeit zu den verwendeten Algorithmen, zur eingesetzten Software und zur einschlägigen Literatur. Dies unterstützte insbesondere die Einrichtung der Simulation, die Entwicklung der Software für die Navigations- und Hinderniserkennung sowie die Strukturierung standardisierter Dokumente wie des Lastenhefts. Zudem halfen seine beruflichen Erfahrungen als Softwaretester dabei, relevante Testfälle zu definieren und zu strukturieren.

Insgesamt bin ich mit der Zusammenarbeit im Projektteam sehr zufrieden und kann mir eine erneute Zusammenarbeit in künftigen Projekten gut vorstellen.